\section{算账：API 报销开支测算}
\label{sec:cost}

本节根据公开 API 价格与典型科研使用强度，测算向中科院科研人员开放 API 报销所需的实际开支。\textbf{所有费用单位为人民币元}。

\medskip
\textbf{基础参数}：
\begin{itemize}
\item 中科院在编科研人员约 \textbf{6 万人}；其中实际会高频使用 AI 工具的约 \textbf{2 万人}（约 1/3）。
\item DeepSeek V4 Pro 国内人民币定价：输入 ¥3 / 百万 token，输出 ¥6 / 百万 token（按原价计算，不计当前 2.5 折优惠）。
\item Claude Opus 4.7 美元定价：输入 \$5 / 百万 token，输出 \$25 / 百万 token；按 1 美元 = 7.2 元换算为输入 \textbf{¥36 / 百万 token}，输出 \textbf{¥180 / 百万 token}。
\item 中文 token 换算系数：\textbf{1 token ≈ 1.3 个汉字}（DeepSeek tokenizer 实测）。
\item 测算\textbf{未计入提示缓存优惠}（缓存命中可使输入费用降至 1/10，实际开支会更低）。
\end{itemize}

\subsection{单次科研用例的费用}

下表给出 4 种典型科研用例的单次调用费用：

\begin{table}[h]
\centering
\renewcommand{\arraystretch}{1.3}
\resizebox{\textwidth}{!}{%
\begin{tabular}{@{}>{\raggedright\arraybackslash}p{5.5cm} >{\centering\arraybackslash}p{2.4cm} >{\centering\arraybackslash}p{2.6cm} >{\centering\arraybackslash}p{2.6cm}@{}}
\toprule
\rowcolor{primary!15}
\textbf{科研用例} & \makecell{\textbf{输入/输出}} & \makecell{\textbf{DeepSeek}} & \makecell{\textbf{Claude}} \\
\midrule
改一段论文摘要 & 2k / 1k & ¥0.012 & ¥0.25 \\
写论文 introduction & 5k / 3k & ¥0.033 & ¥0.72 \\
一篇 arxiv 论文综述 & 40k / 1k & ¥0.13 & ¥1.62 \\
三篇论文交叉对比 & 120k / 3k & ¥0.38 & ¥4.86 \\
\bottomrule
\end{tabular}%
}
\end{table}

\textbf{解读}：DeepSeek 的单次科研调用基本在 \textbf{1 元以内}，即便是最长的多文档对比任务也仅 0.4 元；Claude 的单次费用约为 DeepSeek 的 \textbf{15--20 倍}，但绝对值仍在数元级别——对照单篇论文动辄上万元的版面费、实验耗材、出差差旅等支出，可以忽略不计。

\subsection{单个研究员每月使用强度的费用}

按"每月 22 工作日，每天若干次混合调用"假设，按调用混合比例（50\% 摘要 + 30\% introduction + 15\% 单论文综述 + 5\% 多论文对比）估算每用户月费用：

\begin{table}[h]
\centering
\renewcommand{\arraystretch}{1.3}
\resizebox{\textwidth}{!}{%
\begin{tabular}{@{}>{\raggedright\arraybackslash}p{5.5cm} >{\centering\arraybackslash}p{2.5cm} >{\centering\arraybackslash}p{2.5cm} >{\centering\arraybackslash}p{2.5cm}@{}}
\toprule
\rowcolor{primary!15}
\textbf{用户类型} & \textbf{月调用次数} & \makecell{\textbf{DeepSeek}\\\textbf{\small 月费用}} & \makecell{\textbf{Claude}\\\textbf{\small 月费用}} \\
\midrule
轻度用户（每天 3 次）& 66 & ¥3.5 & ¥55 \\
中度用户（每天 8 次）& 176 & ¥9.5 & ¥146 \\
重度用户（每天 20 次）& 440 & ¥24 & ¥365 \\
\bottomrule
\end{tabular}%
}
\end{table}

\textbf{解读}：即使是\textbf{每天调用 20 次的重度用户}，单月 DeepSeek 费用也仅 ¥24（不到一杯咖啡），单月 Claude 费用约 ¥365。

\subsection{全院级别年度总开支测算}

以\textbf{2 万实际 AI 用户}为基数，按 50\% 轻度 + 35\% 中度 + 15\% 重度的分布估算。下面分两个独立方案：

\medskip
\textbf{方案 A：仅报销 DeepSeek API（国内主力方案）}

\begin{table}[h]
\centering
\renewcommand{\arraystretch}{1.3}
\begin{tabularx}{\textwidth}{@{}>{\raggedright\arraybackslash}X >{\raggedleft\arraybackslash}p{4.5cm}@{}}
\toprule
\rowcolor{primary!15}
\textbf{项目} & \textbf{金额（元）} \\
\midrule
人均月费用（加权平均） & ¥8.6 \\
2 万用户月度总费用 & ¥17.2 万 \\
\rowcolor{primary!8}
\textbf{年度总开支} & \textbf{¥206 万} \\
占中科院年度科研经费比例（按 600 亿计）& \textbf{< 0.004\%} \\
\bottomrule
\end{tabularx}
\end{table}

\medskip
\textbf{方案 B：仅报销 Claude API（高难度任务方案）}

\begin{table}[h]
\centering
\renewcommand{\arraystretch}{1.3}
\begin{tabularx}{\textwidth}{@{}>{\raggedright\arraybackslash}X >{\raggedleft\arraybackslash}p{4.5cm}@{}}
\toprule
\rowcolor{primary!15}
\textbf{项目} & \textbf{金额（元）} \\
\midrule
人均月费用（加权平均） & ¥133 \\
2 万用户月度总费用 & ¥266 万 \\
\rowcolor{primary!8}
\textbf{年度总开支} & \textbf{¥3,200 万} \\
占中科院年度科研经费比例（按 600 亿计）& \textbf{< 0.06\%} \\
\bottomrule
\end{tabularx}
\end{table}

\textbf{两方案对比}：DeepSeek 单走方案年度开支 \textbf{约 200 万}，相当于 \textbf{1 个面上项目的资助额}；Claude 单走方案 \textbf{约 3,200 万}，相当于 \textbf{1 个重点项目的资助额}。\textbf{两者均远低于中科院年度科研经费的千分之一}。

\subsection{对比：自建一套同等性能模型的成本}

为对照参考，下表估算\textbf{从零自建一套对标 DeepSeek V4 Pro / Claude Opus 4.7 性能的大模型}所需投入（业界公开数据估算，不针对任何具体项目）：

\begin{table}[h]
\centering
\renewcommand{\arraystretch}{1.3}
\begin{tabularx}{\textwidth}{@{}>{\raggedright\arraybackslash}X >{\raggedleft\arraybackslash}p{4.5cm}@{}}
\toprule
\rowcolor{primary!15}
\textbf{投入项} & \textbf{金额（元）} \\
\midrule
GPU 集群（1000--2000 张 H100 等效） & ¥2--4 亿 \\
算法团队（30--50 人 × 3 年） & ¥1--1.5 亿 \\
训练数据采购与清洗 & ¥1--2 亿 \\
电费、运维、机房 & ¥3--5 千万 / 年 \\
\midrule
\rowcolor{primary!8}
\textbf{一代模型总投入（3 年摊销）} & \textbf{约 ¥5--8 亿} \\
\textbf{年均摊销} & \textbf{约 ¥1.5--2.5 亿} \\
\bottomrule
\end{tabularx}
\end{table}

\textbf{关键对照}：
\begin{itemize}
\item 即使是最贵的 Claude 全院方案（¥3,200 万 / 年），仅相当于自建模型年均摊销的 \textbf{15\% 左右}。
\item DeepSeek 全院方案（¥206 万 / 年），仅相当于自建模型年均摊销的 \textbf{1\%}。
\item 而且自建模型每代周期为 18--24 个月，下一代发布后需重训，开支不会减少。
\end{itemize}

\subsection{合规与数据安全}

\begin{itemize}
\item \textbf{DeepSeek API}：服务器在境内（北京、杭州），调用与计费均符合国内数据合规要求。可纳入常规科研经费报销，无政策障碍。
\item \textbf{Claude API}：服务器在境外（美国），需通过合规海外节点调用。\textbf{建议仅用于公开数据的科研任务}（如开源论文写作辅助、英文文献处理、开源代码 debug 等），\textbf{涉密、敏感、未发表数据一律不上 Claude}。
\item \textbf{涉密任务}：继续走磐石内网部署或不使用 AI，与本节方案不冲突。
\end{itemize}

\begin{keypoint}
\sffamily\textbf{本节结论}\\[3pt]
\normalfont
\begin{itemize}[leftmargin=18pt,topsep=2pt]
\item \textbf{方案 A}（仅 DeepSeek）：年度开支约 \textbf{¥206 万}，约 \textbf{1 个面上项目}的额度。
\item \textbf{方案 B}（仅 Claude）：年度开支约 \textbf{¥3,200 万}，约 \textbf{1 个重点项目}的额度。
\item 两方案均可与现有磐石部署并行，互不冲突。
\item 与"自建一代同等性能新模型需 ¥5--8 亿"相比，\textbf{API 报销方案的性价比高出 1--2 个数量级}。
\end{itemize}
\end{keypoint}
